% Options for packages loaded elsewhere
% Options for packages loaded elsewhere
\PassOptionsToPackage{unicode}{hyperref}
\PassOptionsToPackage{hyphens}{url}
\PassOptionsToPackage{dvipsnames,svgnames,x11names}{xcolor}
%
\documentclass[
  letterpaper,
  DIV=11,
  numbers=noendperiod]{scrartcl}
\usepackage{xcolor}
\usepackage{amsmath,amssymb}
\setcounter{secnumdepth}{-\maxdimen} % remove section numbering
\usepackage{iftex}
\ifPDFTeX
  \usepackage[T1]{fontenc}
  \usepackage[utf8]{inputenc}
  \usepackage{textcomp} % provide euro and other symbols
\else % if luatex or xetex
  \usepackage{unicode-math} % this also loads fontspec
  \defaultfontfeatures{Scale=MatchLowercase}
  \defaultfontfeatures[\rmfamily]{Ligatures=TeX,Scale=1}
\fi
\usepackage{lmodern}
\ifPDFTeX\else
  % xetex/luatex font selection
\fi
% Use upquote if available, for straight quotes in verbatim environments
\IfFileExists{upquote.sty}{\usepackage{upquote}}{}
\IfFileExists{microtype.sty}{% use microtype if available
  \usepackage[]{microtype}
  \UseMicrotypeSet[protrusion]{basicmath} % disable protrusion for tt fonts
}{}
\makeatletter
\@ifundefined{KOMAClassName}{% if non-KOMA class
  \IfFileExists{parskip.sty}{%
    \usepackage{parskip}
  }{% else
    \setlength{\parindent}{0pt}
    \setlength{\parskip}{6pt plus 2pt minus 1pt}}
}{% if KOMA class
  \KOMAoptions{parskip=half}}
\makeatother
% Make \paragraph and \subparagraph free-standing
\makeatletter
\ifx\paragraph\undefined\else
  \let\oldparagraph\paragraph
  \renewcommand{\paragraph}{
    \@ifstar
      \xxxParagraphStar
      \xxxParagraphNoStar
  }
  \newcommand{\xxxParagraphStar}[1]{\oldparagraph*{#1}\mbox{}}
  \newcommand{\xxxParagraphNoStar}[1]{\oldparagraph{#1}\mbox{}}
\fi
\ifx\subparagraph\undefined\else
  \let\oldsubparagraph\subparagraph
  \renewcommand{\subparagraph}{
    \@ifstar
      \xxxSubParagraphStar
      \xxxSubParagraphNoStar
  }
  \newcommand{\xxxSubParagraphStar}[1]{\oldsubparagraph*{#1}\mbox{}}
  \newcommand{\xxxSubParagraphNoStar}[1]{\oldsubparagraph{#1}\mbox{}}
\fi
\makeatother


\usepackage{longtable,booktabs,array}
\newcounter{none} % for unnumbered tables
\usepackage{calc} % for calculating minipage widths
% Correct order of tables after \paragraph or \subparagraph
\usepackage{etoolbox}
\makeatletter
\patchcmd\longtable{\par}{\if@noskipsec\mbox{}\fi\par}{}{}
\makeatother
% Allow footnotes in longtable head/foot
\IfFileExists{footnotehyper.sty}{\usepackage{footnotehyper}}{\usepackage{footnote}}
\makesavenoteenv{longtable}
\usepackage{graphicx}
\makeatletter
\newsavebox\pandoc@box
\newcommand*\pandocbounded[1]{% scales image to fit in text height/width
  \sbox\pandoc@box{#1}%
  \Gscale@div\@tempa{\textheight}{\dimexpr\ht\pandoc@box+\dp\pandoc@box\relax}%
  \Gscale@div\@tempb{\linewidth}{\wd\pandoc@box}%
  \ifdim\@tempb\p@<\@tempa\p@\let\@tempa\@tempb\fi% select the smaller of both
  \ifdim\@tempa\p@<\p@\scalebox{\@tempa}{\usebox\pandoc@box}%
  \else\usebox{\pandoc@box}%
  \fi%
}
% Set default figure placement to htbp
\def\fps@figure{htbp}
\makeatother





\setlength{\emergencystretch}{3em} % prevent overfull lines

\providecommand{\tightlist}{%
  \setlength{\itemsep}{0pt}\setlength{\parskip}{0pt}}



 


\KOMAoption{captions}{tableheading}
\makeatletter
\@ifpackageloaded{caption}{}{\usepackage{caption}}
\AtBeginDocument{%
\ifdefined\contentsname
  \renewcommand*\contentsname{Table of contents}
\else
  \newcommand\contentsname{Table of contents}
\fi
\ifdefined\listfigurename
  \renewcommand*\listfigurename{List of Figures}
\else
  \newcommand\listfigurename{List of Figures}
\fi
\ifdefined\listtablename
  \renewcommand*\listtablename{List of Tables}
\else
  \newcommand\listtablename{List of Tables}
\fi
\ifdefined\figurename
  \renewcommand*\figurename{Figure}
\else
  \newcommand\figurename{Figure}
\fi
\ifdefined\tablename
  \renewcommand*\tablename{Table}
\else
  \newcommand\tablename{Table}
\fi
}
\@ifpackageloaded{float}{}{\usepackage{float}}
\floatstyle{ruled}
\@ifundefined{c@chapter}{\newfloat{codelisting}{h}{lop}}{\newfloat{codelisting}{h}{lop}[chapter]}
\floatname{codelisting}{Listing}
\newcommand*\listoflistings{\listof{codelisting}{List of Listings}}
\makeatother
\makeatletter
\makeatother
\makeatletter
\@ifpackageloaded{caption}{}{\usepackage{caption}}
\@ifpackageloaded{subcaption}{}{\usepackage{subcaption}}
\makeatother
\usepackage{bookmark}
\IfFileExists{xurl.sty}{\usepackage{xurl}}{} % add URL line breaks if available
\urlstyle{same}
% fallback for those not using the hyperref driver hyperxmp:
\makeatletter
\@ifundefined{xmpquote}{\newcommand{\xmpquote}[1]{#1}}{}
\makeatother
\hypersetup{
  pdftitle={The Logic of Scientific Research},
  pdfauthor={Research Design and Methods in Quantitative Research},
  colorlinks=true,
  linkcolor={blue},
  filecolor={Maroon},
  citecolor={Blue},
  urlcolor={Blue},
  pdfcreator={LaTeX via pandoc}}


\title{The Logic of Scientific Research}
\usepackage{etoolbox}
\makeatletter
\providecommand{\subtitle}[1]{% add subtitle to \maketitle
  \apptocmd{\@title}{\par {\large #1 \par}}{}{}
}
\makeatother
\subtitle{Session 02}
\author{Research Design and Methods in Quantitative Research}
\date{}
\begin{document}
\maketitle

\renewcommand*\contentsname{Table of contents}
{
\setcounter{tocdepth}{1}
\tableofcontents
}

\subsection{\texorpdfstring{The goal of \emph{good}
research}{The goal of good research}}\label{the-goal-of-good-research}

\begin{quote}
To \emph{``produce valid inferences about social and political life''}
(King, Keohane, \& Verba 1994, p.~3)
\end{quote}

. . .

\textbf{What is inference?}

. . .

The process of drawing conclusions from what we \textbf{observe} to what
we do \textbf{not (yet) observe}.

\begin{itemize}
\tightlist
\item
  Quote: KKV p.~3 (ch.~1, ``The Science in Social Science''); basic
  reading today, pp.~3-33
\item
  \textbf{Inference} = using what we observe to learn about what we do
  not observe (KKV p.~8)

  \begin{itemize}
  \tightlist
  \item
    \emph{Descriptive}: from observed facts to unobserved facts (e.g., a
    survey sample -\textgreater{} the population)
  \item
    \emph{Causal}: from data to causal effects (e.g., does canvassing
    raise turnout?)
  \end{itemize}
\item
  Collecting facts alone is not science: the step beyond the data is
  (KKV p.~8)
\item
  Links to session 1 digression (inference vs.~prediction) and to
  sessions 5-7 (descriptive, causal, predictive inference)
\end{itemize}

\subsection{Scientific inference}\label{scientific-inference}

Inferences that draw on evidence and reasoning, following a series of
rules (or methods)

. . .

\emph{The methods make the science, not the content}

\begin{itemize}
\tightlist
\item
  They must be \textbf{public}\\
\item
  They must provide a \textbf{measure of uncertainty}\\
  \emph{(scientific conclusions are always tentative!)}
\end{itemize}

\textbf{KKV's four characteristics of scientific research} (pp.~7-9)

\begin{enumerate}
\def\labelenumi{\arabic{enumi}.}
\tightlist
\item
  \emph{The goal is inference}: go beyond the particular observations
  collected
\item
  \emph{The procedures are public}: explicit, codified methods, so
  others can judge and replicate them (Merton: ``science \ldots{} is
  public, not private'')
\item
  \emph{The conclusions are uncertain}: inference without an uncertainty
  estimate is ``uninterpretable''
\item
  \emph{The content is the method}: science is defined by its rules, not
  its subject matter (Pearson 1892: ``The unity of all science consists
  alone in its method'')
\end{enumerate}

\begin{itemize}
\tightlist
\item
  Implication (p.~9): science is a social enterprise; work does not have
  to be beyond criticism to contribute
\item
  Slide shows 2-4 as bullets; 1 was the previous slide
\end{itemize}

\subsection{}\label{section}

\textbf{Share your answers!}

{\emph{Seminar questions 2}}

\begin{itemize}
\item
  Which of KKV's four characteristics did you pick, and why does it
  matter?
\item
  Do qualitative and quantitative research share the same logic of
  inference?
\end{itemize}

\begin{center}
\includegraphics[width=0.75\linewidth,height=\textheight,keepaspectratio]{volunteers.png}
\end{center}

\begin{itemize}
\tightlist
\item
  Picks up seminar questions 2 (Q2 and Q3); answers were due Wednesday
  30.09, 12:00
\item
  Q2: ask 2-3 students which characteristic they picked and their
  example; map each answer onto the list (previous notes)
\item
  Q3 (one logic of inference): KKV's claim (pp.~3-4) is that qualitative
  and quantitative research differ in style and technique, not in the
  underlying logic

  \begin{itemize}
  \tightlist
  \item
    Counterpoint if the debate takes off: Brady \& Collier (eds.),
    \emph{Rethinking Social Inquiry} (2004), argue KKV judge qualitative
    work by quantitative standards
  \item
    Our course: quantitative focus, but small-N designs in session 10
  \end{itemize}
\item
  Keep it short (\textasciitilde5 min); the RQ exercise comes later
\end{itemize}

\subsection{Components of a research
design}\label{components-of-a-research-design}

\begin{itemize}
\tightlist
\item
  \textbf{Research question} (this session)\\
\item
  \textbf{Theory} (session 3)\\
\item
  \textbf{Data} (session 4)\\
\item
  \textbf{Methods}: depend on inferential goals (sessions 5, 6, 7),
  resources, and data availability (sessions 8, 9, 10).\\
  We will focus on \textbf{quantitative methods} (sessions 11, 12, 13)
\end{itemize}

\begin{itemize}
\tightlist
\item
  KKV p.~13: four components of every research design: the research
  question, the theory, the data, and the use of the data
\item
  They are rarely developed in this order; researchers go back and forth
  (KKV pp.~12-13: ``taking down the scaffolding'')
\item
  Course map (syllabus schedule): 3 theory - 4 data and measurement -
  5/6/7 descriptive, causal, predictive inference - 8/9/10 experimental,
  large-N, small-N - 11 statistical testing - 12/13 regression - 14
  conclusion
\item
  Today and session 3 happen on the same afternoon: research question
  now, theory at 16:15
\end{itemize}

\subsection{The research question}\label{the-research-question}

There is no logical recipe for having good ideas.

. . .

It depends on interest, creativity, knowledge\ldots{} and luck.

. . .

However, some guidelines:

\begin{itemize}
\tightlist
\item
  Must be \textbf{relevant}\\
\item
  Answering it should make a \textbf{contribution}\\
  (know your literature! or goal --- e.g., private company)\\
\item
  Should be \textbf{researchable} (realistic plans to answer it!)
\item
  Should be \textbf{focused}
\end{itemize}

\begin{itemize}
\tightlist
\item
  ``No logical recipe'': KKV p.~14, quoting Popper: ``there is no such
  thing as a logical method of having new ideas''
\item
  KKV's two criteria (p.~15), also in seminar question 1:

  \begin{enumerate}
  \def\labelenumi{\arabic{enumi}.}
  \tightlist
  \item
    \emph{Important in the real world}: consequential for political,
    social or economic life
  \item
    \emph{A specific contribution to an identifiable scholarly
    literature}
  \end{enumerate}
\item
  Six ways to contribute (KKV pp.~16-17): (1) test an important
  hypothesis nobody has studied systematically; (2) challenge an
  accepted hypothesis; (3) help resolve a controversy; (4) examine
  unquestioned assumptions; (5) show an important topic has been
  overlooked; (6) apply theories or evidence from one literature to
  another
\item
  Personal motivation is fine but not a justification: ``no one cares
  what we think - the scholarly community only cares what we can
  demonstrate'' (p.~15)
\item
  \emph{Researchable}: KKV p.~18: a topic that cannot be refined into a
  specific project permitting valid inference ``should be modified along
  the way or abandoned''
\item
  \emph{Focused} is our addition: narrow enough to answer with one
  design
\item
  Private-sector aside: there the ``contribution'' is to the
  organisation's goal, not a literature
\end{itemize}

\subsection{Types of research
questions}\label{types-of-research-questions}

{\def\LTcaptype{none} % do not increment counter
\begin{longtable}[]{@{}
  >{\raggedright\arraybackslash}p{(\linewidth - 4\tabcolsep) * \real{0.2206}}
  >{\raggedright\arraybackslash}p{(\linewidth - 4\tabcolsep) * \real{0.1912}}
  >{\raggedright\arraybackslash}p{(\linewidth - 4\tabcolsep) * \real{0.5882}}@{}}
\toprule\noalign{}
\begin{minipage}[b]{\linewidth}\raggedright
Question
\end{minipage} & \begin{minipage}[b]{\linewidth}\raggedright
Focus
\end{minipage} & \begin{minipage}[b]{\linewidth}\raggedright
Example
\end{minipage} \\
\midrule\noalign{}
\endhead
\bottomrule\noalign{}
\endlastfoot
\textbf{What} & Descriptive & What is voter turnout in country X? \\
\textbf{How} & Explanatory & How do campaigns influence turnout? \\
\textbf{Why} & Causal & Why do some groups vote more than others? \\
\end{longtable}
}

\begin{itemize}
\tightlist
\item
  What/How/Why maps roughly onto the inferential goals of sessions 5-7:
  descriptive (what), mechanisms (how), causal (why)
\item
  Descriptive questions are legitimate research questions (session 5
  readings: Gerring 2012, \emph{Mere Description}; de Kadt \&
  Grzymala-Busse 2025)
\item
  Ask: which type is your seminar-question RQ?
\end{itemize}

\subsection{Group exercise: Assess your research
question!}\label{group-exercise-assess-your-research-question}

. . .

\begin{itemize}
\tightlist
\item
  Take the research question you wrote for the \textbf{seminar
  questions} (question 1)\\
\item
  Choose a partner next to you and share your research question\\
\item
  Your partner should do the same with you\\
\item
  Assess your partner's question (\textbf{three minutes}):

  \begin{itemize}
  \tightlist
  \item
    Does it meet the criteria above?\\
  \item
    How could it be improved?
  \end{itemize}
\end{itemize}

\begin{itemize}
\tightlist
\item
  Students use the RQ from seminar question 1 (uploaded by Wednesday
  12:00); no new writing step
\item
  \textasciitilde5 minutes in pairs; walk around, note 2-3 good examples
  to call on
\item
  Criteria to check: relevant, contribution, researchable, focused
  (previous slides); KKV's two criteria are in the seminar question
\item
  Anyone without a question: pair them with someone and let them assess
\end{itemize}

\subsection{}\label{section-1}

\textbf{Share your research question!}

{\emph{Seminar questions 2}}

\begin{itemize}
\item
  What is your question?
\item
  Is it relevant, researchable and focused?
\item
  How did your partner improve it?
\end{itemize}

\begin{center}
\includegraphics[width=0.75\linewidth,height=\textheight,keepaspectratio]{volunteer_meme.png}
\end{center}

\begin{itemize}
\tightlist
\item
  2-3 volunteers; for each: which criterion is weakest, and how did the
  partner fix it?
\item
  Typical problems: too broad (``What drives populism?''), normative
  (``Should we ban \ldots?''), not researchable (no data), or already
  answered
\item
  Keep the questions in mind: session 3 builds on them (seminar question
  3 asks for a hypothesis derived from the same RQ)
\end{itemize}

\subsection{Conclusion}\label{conclusion}

A good \textbf{research question} is the first component of an effective
research design

. . .

Later today (session 3, 16:15), we will examine more of your research
questions and assess how to connect them with a \textbf{theory and
hypotheses}

\begin{itemize}
\tightlist
\item
  The RQ is component 1 of 4 (KKV p.~13)
\item
  Session 3 today at 16:15: from RQ to theory, models and hypotheses
  (seminar question 3 builds on seminar question 1 of today)
\end{itemize}

\begin{center}\rule{0.5\linewidth}{0.5pt}\end{center}

The ultimate goal of this course is to \textbf{\emph{think like a social
scientist}}

. . .

Connect real-world observations to theory and methods

. . .

\textbf{Remain skeptical!}

\begin{itemize}
\tightlist
\item
  ``Thinking like a social scientist'': KKV section 1.3.4, pp.~32-33,
  ``Skepticism and Rival Hypotheses''
\item
  KKV example: Japanese eat less red meat and have fewer heart attacks
  than Americans. A skeptic asks: are the data accurate and comparable?
  What else differs (other diet, genes, lifestyle)?
\item
  Ask these questions of your own research, not just of others'
\item
  Bridge: that is exactly what the presentations do with each applied
  reading
\end{itemize}

\subsection{Up next: applied reading}\label{up-next-applied-reading}

Presentation of ongoing research by the lecturer, in the \textbf{same
format} as your presentations:

\begin{itemize}
\tightlist
\item
  Klüser, S., Klüser, K. J., \& Hoes, E. (2026). AI-Generated Explicit
  Deepfakes Damage Politicians' Perceived Leadership Competence,
  Trustworthiness and Electoral Prospects. \emph{Revised and
  resubmitted, Communications Psychology}
\end{itemize}

Part I: summary and critique (\textasciitilde10 min) · Part II:
alternative design (\textasciitilde15 min) · discussion

\begin{itemize}
\tightlist
\item
  Switch to the PowerPoint deck:
  02\_session02/presentation\_slides\_session02\_KLUESER.pptx (speaker
  notes inside)
\item
  Same format and rubric as the student presentations: \textasciitilde10
  min summary and critique, \textasciitilde15 min alternative design,
  then discussion (syllabus section 5, grading sheet)
\item
  The paper is the applied reading in OLAT (July 2026 revision, R\&R at
  Communications Psychology)
\item
  Content note: the treatment videos in the deck are blurred; no
  explicit imagery is shown
\end{itemize}

\subsection{Next session (01.10.26 - 16:15, room
2.A45)}\label{next-session-01.10.26---1615-room-2.a45}

\textbf{\emph{Theory and research design}} · Basic reading: KW, ch.~2
(pp.~25--52)

. . .

Presentation by \textbf{Alessio Gallonetto, Lea Hafen and Benjamin
Seiwert}:\\
Turnbull-Dugarte \& Wagner (2025), \emph{Heroes and villains: motivated
projection of political identities}

\begin{itemize}
\tightlist
\item
  Same day: session 3, 16:15-18:00, room 2.A45 (syllabus schedule)
\item
  Basic reading: KW (Kellstedt \& Whitten), ch.~2, ``The Art of Theory
  Building'', pp.~25-52
\item
  First student presentation: Alessio Gallonetto, Lea Hafen, Benjamin
  Seiwert (three presenters in sessions 3 and 4)
\item
  This group may upload their slides on the morning of 1 October instead
  of Wednesday noon (OLAT message of 23.09)
\item
  Seminar questions 3 were due Wednesday 30.09, 12:00, together with
  today's
\end{itemize}

\begin{center}\rule{0.5\linewidth}{0.5pt}\end{center}

\begin{center}
\pandocbounded{\includegraphics[keepaspectratio]{questions.png}}
\end{center}

\begin{itemize}
\tightlist
\item
  Open floor: questions on KKV, the research question exercise, or the
  presentation format
\end{itemize}

\subsection{\texorpdfstring{Thanks! {🙂}}{Thanks! 🙂}}\label{thanks}

\href{mailto:klueser@ipz.uzh.ch}{\nolinkurl{klueser@ipz.uzh.ch}}

\begin{itemize}
\tightlist
\item
  Contact: klueser@ipz.uzh.ch; meetings via Zoom on request within a
  week (syllabus: no fixed office hours)
\end{itemize}




\end{document}
